%% --------------------------------------------------------------------
%% Template UTEC Tesis
%% --------------------------------------------------------------------
% Este es el archivo que se debe compilar.
% 
% Este template ha sido modificado y actualizado por Eduardo Castro y Roosevelt Ubaldo en base a lo trabajado por Víctor Murray, Oscar Ramos y Juan Carlos Barbaran.
%
% Última actualización: Mayo, 2022

\documentclass[a4paper,12pt,oneside]{tesisutec}

\selectlanguage{spanish}

%% Paquetes
\usepackage[utf8]{inputenc}
\usepackage[square, numbers, comma, sort&compress]{natbib}
% Libreria de idioma
\usepackage[spanish]{babel}
% Libreria para posicionamiento
\usepackage{float}
% Librerias para insertar codigos
\usepackage[spanish,onelanguage,ruled,vlined]{algorithm2e}
\usepackage{verbatim} 
% Librería para hipervínculos
\usepackage{hyperref}
 % Librería necesaria para arreglar el orden de referencias en overleaf.com
\usepackage{notoccite}

% Incluir acá los paquetes adicionales que deseas. 
% Ubicación de los imágenes.
\graphicspath{images/}

\makeatletter
% Reinsert missing \algbackskip
\def\algbackskip{\hskip-\ALG@thistlm}
\makeatother

\hypersetup{urlcolor=blue, colorlinks=true}

\begin{document}

\frontmatter

\department{}
\degree{bachiller }
\major{Ingeniería Mecatrónica}

\title{Monitoreo de granjas de jaula para salmonicultura usando visión computacional en un ROV}

\author{Jim Elashir Fabián Gonzales \orcid{0000-0000-0000-0000}} % Es obligatorio agregar ORCID del alumno
\supervisor{Oscar E. Ramos \orcid{0000-0000-0000-0000}} % Es obligatorio agregar ORCID del asesor

\date{2023}

\maketitle

\setstretch{1.5}


% %% ============================================================================
%  Dedicatoria
%% ============================================================================

\dedicatory{
xxxxxxxxxxxxxxxxxxxxxxxxxxxxxxxxxxxxxxxxxxxxxxxxxxx
xxxxxxxxxxxxxxxxxxxxxxxxxxxxxxxxxxxxxxxxxxxxxxxxxxx
xxxxxxxxxxxxxxxxxxxxxxxxxxxxxxxxxxxxxxxxxxxxxxxxxxx
xxxxxxxxxxxxxxxxxxxxxxxxxxxxxxxxxxxxxxxxxxxxxxxxxxx
xxxxxxxxxxxxxxxxxxxxxxxxxxxxxxxxxxxxxxxxxxxxxxxxxxx
XXxxxxxxxxxxxxxxxxxxxxxxxxxxxxxxxxxxxxxxxxxxxxxxxx
}

% %% ============================================================================
%  Pagina de Agradecimientos
%% ============================================================================

\acknowledgements{

xxxxxxxxxxxxxxxxxxxxxxxxxxxxxxxxxxxxxxxxxxxxxxxxxxx
xxxxxxxxxxxxxxxxxxxxxxxxxxxxxxxxxxxxxxxxxxxxxxxxxxx
xxxxxxxxxxxxxxxxxxxxxxxxxxxxxxxxxxxxxxxxxxxxxxxxxxx
xxxxxxxxxxxxxxxxxxxxxxxxxxxxxxxxxxxxxxxxxxxxxxxxxxx
xxxxxxxxxxxxxxxxxxxxxxxxxxxxxxxxxxxxxxxxxxxxxxxxxxx
xxxxxxxxxxxxxxxxxxxxxxxxxxxxxxxxxxxxxxxxxxxxxxxxxxx

}



\tableofcontents

\newpage

% \listoftables

\newpage

% \listoffigures

\addtocontents{toc}{\vspace{1.5em}}

%% ============================================================================
\mainmatter
\pagestyle{fancy}

% \customchapter{RESUMEN}

Lorem ipsum dolor sit amet, cum tollit aeterno ut. Sea amet iuvaret placerat in, saperet impedit consulatu mei ne, quo no alii vituperatoribus. No mutat dolorum habemus pro, vis rebum accusam no, meis partiendo necessitatibus cu sit. Zril mediocrem scriptorem cu quo, ex duo mucius doming legendos. Posse nemore cu cum. Et facer pertinax sed, atqui ignota antiopam in pro, dolore oportere ius ex.

An justo nonumy mnesarchum quo, eruditi cotidieque in sed. Consulatu urbanitas conceptam eam id, liber accusamus eu mel. Ius iusto albucius elaboraret id, dico partiendo contentiones nec in, sint eruditi per ne. No eruditi probatus cum, no mei sumo facer torquatos. Sea dicunt facilisis corrumpit at, audiam accusam eligendi qui ne.

Facete officiis invenire sit ex, facete causae fabellas ei qui. Et pertinacia contentiones definitiones nam. Cum moderatius signiferumque in, blandit senserit cum no. Mea appareat posidonium contentiones an, no case accommodare cum, ei sea oporteat percipitur. 
Ea congue necessitatibus sed, no adhuc graeco usu. Eam harum equidem fierent ne, duo id facer consequuntur signiferumque. 
Dolore molestie phaedrum nec ei, ei eos noluisse constituto. Eruditi propriae in mea, vero laboramus cotidieque ad duo. Tacimates constituto est ex, ad mei clita sapientem evertitur. \\


\noindent \textbf{Palabras clave:}\\
\noindent Palabra clave 1; Palabra clave 2; Palabra clave 3; Palabra clave 4
% \customchapter{ABSTRACT} 

\begin{center}
\large \vspace{-1.5cm} \textbf{TÍTULO DE LA TESIS EN INGLÉS}
\end{center}

Lorem ipsum dolor sit amet, with the eternal God, that takes away the. Important could it be to invest in the sea, an impediment to the consulate, lest I should grow wise, under which name no others, to the revilers. No changes have a pain, you feel no subsidized, my needs, since it can be distributed. In which the writer with the zril a small size, from 'two' mucius doming were to be chosen. A man can do with him. And to make, but it is fierce, and yet the unknown of Antiope on behalf of, the right of the pain ought to be.

Is this just mnesarchum running smoothly, but instructed the daily. Consulate refinement to this concept, a book accusing football honey. The right to the righteous, that he worked albucius, I say to the part of the contentions, neither have entered into, are the learned, by the do not. No approved as instructed, no curling but do take me. Antioxidants dough at sea say, I feel they should be chosen.

Elegantly offices to find out, salty stories to the cause. For definitions and resolution of conflicts. When measured signiferumque to sit when he feels no. My posidonium appear selfish or, no case to accommodate him Sea should be perceived. Planning needs but that no further Greek practice. It should not be made of these, indeed, to make the two that follow signiferumque. \\


\noindent \textbf{Keywords:}\\
\noindent Keyword 1; Keyword 2; Keyword 3; Keyword 4


 
\customchapter{INTRODUCCIÓN} 

\introsection{Presentación del tema de investigación}

Lorem ipsum dolor sit amet, consectetur adipiscing elit. In nec massa venenatis, sodales magna id, tempus lacus. Nullam nec dolor eu ex luctus molestie. Ut quam libero, tristique et efficitur at, consequat vel velit. Aliquam non sem eu libero aliquet condimentum. Curabitur lobortis neque purus, vel pellentesque lorem tincidunt a. Pellentesque vitae nunc luctus, rutrum sapien sit amet, cursus eros. In tempus, nulla sit amet pharetra iaculis, urna odio dapibus quam, sed commodo magna massa sit amet sem. Proin id finibus nulla. Vestibulum ante ipsum primis in faucibus orci luctus et ultrices posuere cubilia curae.

\introsection{ Descripción de la situación problemática}

Lorem ipsum dolor sit amet, consectetur adipiscing elit. In nec massa venenatis, sodales magna id, tempus lacus. Nullam nec dolor eu ex luctus molestie. Ut quam libero, tristique et efficitur at, consequat vel velit. Aliquam non sem eu libero aliquet condimentum. Curabitur lobortis neque purus, vel pellentesque lorem tincidunt a. Pellentesque vitae nunc luctus, rutrum sapien sit amet, cursus eros. In tempus, nulla sit amet pharetra iaculis, urna odio dapibus quam, sed commodo magna massa sit amet sem. Proin id finibus nulla. Vestibulum ante ipsum primis in faucibus orci luctus et ultrices posuere cubilia curae.

\introsection{Formulación del problema}

Lorem ipsum dolor sit amet, consectetur adipiscing elit. In nec massa venenatis, sodales magna id, tempus lacus. Nullam nec dolor eu ex luctus molestie. Ut quam libero, tristique et efficitur at, consequat vel velit. Aliquam non sem eu libero aliquet condimentum. Curabitur lobortis neque purus, vel pellentesque lorem tincidunt a. Pellentesque vitae nunc luctus, rutrum sapien sit amet, cursus eros. In tempus, nulla sit amet pharetra iaculis, urna odio dapibus quam, sed commodo magna massa sit amet sem. Proin id finibus nulla. Vestibulum ante ipsum primis in faucibus orci luctus et ultrices posuere cubilia curae.

\introsection{Objetivos de investigación}

Lorem ipsum dolor sit amet, consectetur adipiscing elit. In nec massa venenatis, sodales magna id, tempus lacus. Nullam nec dolor eu ex luctus molestie. Ut quam libero, tristique et efficitur at, consequat vel velit. Aliquam non sem eu libero aliquet condimentum. Curabitur lobortis neque purus, vel pellentesque lorem tincidunt a. Pellentesque vitae nunc luctus, rutrum sapien sit amet, cursus eros. In tempus, nulla sit amet pharetra iaculis, urna odio dapibus quam, sed commodo magna massa sit amet sem. Proin id finibus nulla. Vestibulum ante ipsum primis in faucibus orci luctus et ultrices posuere cubilia curae.

\introsection{Justificación}

Lorem ipsum dolor sit amet, consectetur adipiscing elit. In nec massa venenatis, sodales magna id, tempus lacus. Nullam nec dolor eu ex luctus molestie. Ut quam libero, tristique et efficitur at, consequat vel velit. Aliquam non sem eu libero aliquet condimentum. Curabitur lobortis neque purus, vel pellentesque lorem tincidunt a. Pellentesque vitae nunc luctus, rutrum sapien sit amet, cursus eros. In tempus, nulla sit amet pharetra iaculis, urna odio dapibus quam, sed commodo magna massa sit amet sem. Proin id finibus nulla. Vestibulum ante ipsum primis in faucibus orci luctus et ultrices posuere cubilia curae.

\introsection{Alcance y limitaciones / restricciones} %[opcional]

Lorem ipsum dolor sit amet, consectetur adipiscing elit. In nec massa venenatis, sodales magna id, tempus lacus. Nullam nec dolor eu ex luctus molestie. Ut quam libero, tristique et efficitur at, consequat vel velit. Aliquam non sem eu libero aliquet condimentum. Curabitur lobortis neque purus, vel pellentesque lorem tincidunt a. Pellentesque vitae nunc luctus, rutrum sapien sit amet, cursus eros. In tempus, nulla sit amet pharetra iaculis, urna odio dapibus quam, sed commodo magna massa sit amet sem. Proin id finibus nulla. Vestibulum ante ipsum primis in faucibus orci luctus et ultrices posuere cubilia curae.
\chapter{REVISIÓN CRÍTICA DE LA LITERATURA}

\section{Robótica submarina en acuicultura.}

holi \cite{Reumann2012}.

\section{Segundo subtítulo}

\subsection{División del segundo subtítulo}

La industria acuícola 
%
\begin{equation}
  \label{eq:1}
  a+b=\sqrt{\frac{4}{3}},
\end{equation}
%
donde $a$ y $b$ son escalares.

Entre las causas se pueden considerar la falta de tiempo asignado al taller,
las limitaciones en cuanto a laboratorios, el poco personal, la falta de una
correcta selección de los contenidos, así como, una calificación que asegure el
cumplimiento de los objetivos. Un ejemplo de cómo poner una figura se muestra en la Fig. \ref{fig:diagram2}, y es importante recordar la relación mostrada en \eqref{eq:1}.


\begin{figure} 
\begin{center}
\begin{tabular}{cc}
\includegraphics[height=3cm]{images/logo_utec.png} &
\includegraphics[height=2cm]{images/logo_utec.png} \\
(a) & (b)
\end{tabular}
\caption{\label{fig:diagram2}Scheme showing the architecture of a generic kinematic task. (a) Logo en tamaño de 3 centímetros. (b) Logo en tamaño de 2 centímetros.}
\end{center}
\end{figure}

\section{Tercer subtítulo}

\subsection{División del tercer subtítulo}

En la actualidad, las unidades de información hacen frente a muchos cambios
debido a los avances tecnológicos, explosión informativa, nuevos recursos y
soportes, por lo cual se implementan servicios innovadores que les permitan a
sus usuarios tener acceso a muchas fuentes de información. Para asegurar el
acceso y uso de los servicios los usuarios requieren poseer una serie de
habilidades que les permitan identificar, recuperar, manejar, discernir,
organizar, utilizar y comunicar la información de manera eficaz para la toma de
decisiones.

Entre las causas se pueden considerar la falta de tiempo asignado al taller,
las limitaciones en cuanto a laboratorios, el poco personal, la falta de una
correcta selección de los contenidos, así como, una calificación que asegure el
cumplimiento de los objetivos.

\begin{table}[H]
    \centering
    \begin{tabular}{c|c}
        Tiempo($s$) & Distancia($m$) \\
        \hline
        10 & 23 \\
        20 & 33 \\
        30 & 43 \\
        40 & 53 \\
    \end{tabular}
    \caption{Tiempos versus distancia.}
    \label{tab:tiempo_versus_distancia}
\end{table}
\chapter{MARCO TEÓRICO}

Inicie aquí el texto, utilizando sangría de 1.25 cm. en el primer
párrafo. Continúe el segundo párrafo.

Continúe el segundo párrafo
\section{Vehículo operado remotamente - ROV}
\section{Sensores exteroceptivos}

\section{Visión Computacional}
\subsection{Visual SLAM}
\subsubsection{Partes }
\subsubsection{Inertial Visual SLAM}
En la actualidad, las unidades de información hacen frente a muchos cambios
debido a los avances tecnológicos, explosión informativa, nuevos recursos y
soportes, por lo cual se implementan servicios innovadores que les permitan a
sus usuarios tener acceso a muchas fuentes de información. Para asegurar el
acceso y uso de los servicios los usuarios requieren poseer una serie de
habilidades que les permitan identificar, recuperar, manejar, discernir,
organizar, utilizar y comunicar la información de manera eficaz para la toma de
decisiones.

Entre las causas se pueden considerar la falta de tiempo asignado al taller,
las limitaciones en cuanto a laboratorios, el poco personal, la falta de una
correcta selección de los contenidos, así como, una calificación que asegure el
cumplimiento de los objetivos.


% \section{Segundo subtítulo}

En la actualidad, las unidades de información hacen frente a muchos cambios
debido a los avances tecnológicos, explosión informativa, nuevos recursos y
soportes, por lo cual se implementan servicios innovadores que les permitan a
sus usuarios tener acceso a muchas fuentes de información. Para asegurar el
acceso y uso de los servicios los usuarios requieren poseer una serie de
habilidades que les permitan identificar, recuperar, manejar, discernir,
organizar, utilizar y comunicar la información de manera eficaz para la toma de
decisiones.

Entre las causas se pueden considerar la falta de tiempo asignado al taller,
las limitaciones en cuanto a laboratorios, el poco personal, la falta de una
correcta selección de los contenidos, así como, una calificación que asegure el
cumplimiento de los objetivos.


% \chapter{MARCO REFERENCIAL (LEGAL, GEOGRÁFICO, HISTÓRICO) (opcional)}

Inicie aquí el texto, utilizando sangría de 1.25 cm. en el primer
párrafo. Continúe el segundo párrafo.

Continúe el segundo párrafo

\section{Primer subtítulo}

En la actualidad, las unidades de información hacen frente a muchos cambios
debido a los avances tecnológicos, explosión informativa, nuevos recursos y
soportes, por lo cual se implementan servicios innovadores que les permitan a
sus usuarios tener acceso a muchas fuentes de información. Para asegurar el
acceso y uso de los servicios los usuarios requieren poseer una serie de
habilidades que les permitan identificar, recuperar, manejar, discernir,
organizar, utilizar y comunicar la información de manera eficaz para la toma de
decisiones.

Entre las causas se pueden considerar la falta de tiempo asignado al taller,
las limitaciones en cuanto a laboratorios, el poco personal, la falta de una
correcta selección de los contenidos, así como, una calificación que asegure el
cumplimiento de los objetivos.


% \chapter{MARCO METODOLÓGICO} \label{MARCO_METODOLOGICO}

Esta investigación práctica se basa en la implementación de un algoritmo que utiliza técnicas de Deep Learning, SLAM visual-inercial-presión y procesamiento de imágenes para monitorear las redes de las granjas de jaulas de salmón. A continuación, se detalla la metodología que se seguirá para cada uno de estos componentes. 

\section{Desarrollo del algoritmo de Deep Learning}

El primer paso en el desarrollo del algoritmo de Deep Learning será la recopilación de un conjunto de datos de imágenes de redes de granjas de jaulas de salmón para crear el dataset. Estas imágenes serán etiquetadas manualmente para indicar la ubicación de los huecos en la red. Este conjunto de datos se dividirá en un conjunto de entrenamiento y un conjunto de prueba.

Para el entrenamiento del modelo de Deep Learning, se utilizará una arquitectura de red neuronal convolucional (CNN) basada en Mask CRNN presentada en \cite{cite:Zhang}. Esta arquitectura ha demostrado ser eficaz en tareas de reconocimiento de imágenes y se espera que sea capaz de identificar los huecos en las redes de las granjas de jaulas. Durante el entrenamiento, se utilizarán técnicas de aumento de datos para mejorar la robustez del modelo y evitar el sobre ajuste. 

Una vez que el modelo haya sido entrenado y validado, se integrará en un robot submarino. Para ello, se utilizará un sistema operativo robótico (ROS 2) que permitirá la comunicación entre el modelo de Deep Learning y el hardware del robot. Este sistema permitirá que el modelo de Deep Learning procese las imágenes capturadas por una cámara en tiempo real y detecte los huecos en la red.

\section{Aplicación de técnicas de procesamiento de imágenes}

Para la identificación de áreas sucias en la red, se utilizarán técnicas de procesamiento de imágenes presentadas en \cite{cite:Zhao}. En primer lugar, se utilizará el método de Otsu para convertir las imágenes en escala de grises a imágenes binarias. Este método selecciona un umbral que minimiza la varianza intra clase en la imagen, lo que resulta en una segmentación efectiva de las áreas sucias y limpias de la red.

Una vez que se ha obtenido la imagen binaria, se utilizará el Histograma de Gradientes Orientados (HOG) para extraer características de la imagen. El HOG es una técnica que cuenta las ocurrencias de las orientaciones de los gradientes en porciones de la imagen. Estas características serán útiles para identificar las áreas sucias de la red, ya que las áreas sucias y limpias tendrán diferentes distribuciones de orientaciones de gradientes.

Finalmente, se calculará el porcentaje de suciedad en la red contando el número de píxeles en las áreas identificadas como sucias y dividiéndolo por el número total de píxeles en la imagen. Este porcentaje de suciedad será una medida cuantitativa de la limpieza de la red.

\section{Aplicación de SLAM visual-inercial-presión} \label{VIP_SLAM}

La realizará una simulación e implementación del SLAM visual-inercial-presión presentada en \cite{cite:Ferrera}, el cual se realizará en tres etapas. En la primera etapa, se implementará un SLAM visual utilizando una cámara monocular. Se utilizará el algoritmo de SLAM visual ORB-SLAM2 \cite{murTRO2015} y se adaptará para trabajar con las imágenes de una cámara Kinect junto a un turtlebot Waffle pi. La simulación se realizará en ROS Noetic con Gazebo y con RViz la reconstrucción. Este algoritmo permitirá construir un mapa 3D de su entorno y localizarse dentro de este mapa. 

Cabe resaltar que en \cite{cite:Ferrera} se implementó UW-VO SLAM. Sin embargo, no se ha encontrado un repositorio de código abierto. Por tal motivo, dentro de la comparación que se realiza junto a otros métodos de monocular SLAM, ORB-SLAM es el segundo con menor error. 

En la segunda etapa, se añadirá la información de presión al algoritmo de SLAM visual. Para ello, se utilizará un sensor de presión que se incorporará en un sistema mecatrónico junto a la cámara. La información de presión se utilizará para mejorar la estimación de la profundidad en el algoritmo de SLAM. Esto permitirá conocer su profundidad en el agua, lo que es crucial para la navegación submarina.

En la tercera etapa, se añadirá la información inercial al algoritmo de SLAM. Para ello, se utilizará un IMU que se instalará al sistema mecatrónico. La información inercial se utilizará para mejorar la estimación de la orientación y la velocidad en el algoritmo de SLAM. Esto permitirá conocer su orientación y velocidad del sistema.

% \customchapter{CONCLUSIONES} 

Lorem ipsum dolor sit amet, dicant volutpat vulputate sea at. Et nam volumus voluptua, ea noster omnesque has. Qui ut graeci melius platonem. Sit fugit prompta dignissim cu. Soluta meliore repudiare vis in, eum id liber maiestatis contentiones.

Ea sit suas melius verear. Fugit facilisis at vel, apeirian appellantur ex nec. Vis id quis elit qualisque. Ludus sadipscing eu eum.

Ea inani insolens dissentiunt eos, wisi efficiendi ex vim, modus justo inani et pro. Eam at modus definiebas, id eos philosophia disputationi, eu qui facete habemus. In dico unum neglegentur mea, ne phaedrum disputationi his. Tibique deseruisse mel ne, latine aperiri phaedrum duo ex. Erant corrumpit ut duo, in erant possim reprehendunt nec. Aeque nonumes minimum te eam, possim constituto qui ei, in has suas quidam molestie.

Vel no intellegat voluptaria, ei sed causae scaevola comprehensam. No erant cetero dissentiet per, mei an nulla scaevola. Euismod delicata vis eu, in odio docendi persecuti sea. Eos fabellas accusata at. His eu reque altera sadipscing, id ius regione lobortis expetenda, minim timeam dignissim vix eu.

Unum duis definitiones vel in. Nulla dicit mei ea, odio nulla an pri. Cu docendi torquatos usu. Possit maiorum rationibus ea duo, eu vitae primis postulant sed, mel erat officiis principes an. Pro saepe regione volumus id, debet reformidans ex his. Sale altera moderatius duo ea, an fastidii volutpat per, et mel maiorum corpora.

Modo malis et eam. Veniam equidem mei ad, mucius facilisi et pro. Postea eripuit neglegentur vis no. Id vis quis eleifend appellantur, sit integre singulis ei.
% \chapter*{\center \Large RECOMENDACIONES} 
\addcontentsline{toc}{section}{\bfseries RECOMENDACIONES} 
\markboth{RECOMENDACIONES}{RECOMENDACIONES} 

Lorem ipsum dolor sit amet, mei no laboramus gloriatur, no timeam euripidis qui, eu eos aperiam patrioque accommodare. Mei probo accommodare an, ex case minim salutatus sit. Ad doming impetus mei, placerat verterem has no. His assum quaeque et, ea aliquid dissentiet eum. Copiosae evertitur cum te, sea clita disputationi ea.

Libris ridens malorum cu vis. Per vitae eirmod cotidieque cu. Cum unum mucius alterum eu, nam enim regione appellantur id. Cum putant dignissim te, est dolore essent definiebas ei. Facilisi convenire vix in, dictas delenit et per. Pro tota affert argumentum cu. % sección opcional

%% ============================================================================
\renewcommand{\bibname}{\hfill\Large\bf{REFERENCIAS BIBLIOGRÁFICAS}\hfill}

\bibliographystyle{IEEEtran} % Estableciendo el estilo de citas IEEE tram.

\bibliography{referencias} % Recibe las referencias de IEEE

% \chapter*{\center \Large ANEXOS} 
\addcontentsline{toc}{section}{\bfseries ANEXOS} 
\markboth{ANEXOS}{ANEXOS} 

\par Los algoritmos desarrollados .....
    

\end{document}
